\documentclass[11pt]{article}

\usepackage[margin=0.65in]{geometry}
\usepackage{xcolor}
\usepackage{tcolorbox}
\usepackage{tabularx}
\usepackage{array}
\usepackage{enumitem}
\usepackage{amsmath,amssymb}
\usepackage{titlesec}
\usepackage{fancyhdr}
\usepackage{tikz}
\usepackage{hyperref}

\definecolor{NovelPurple}{HTML}{4B3FA7}
\definecolor{LightPurple}{HTML}{F2EFFF}
\definecolor{SoftGray}{HTML}{F7F7F7}

\newcommand{\circnum}[1]{%
\tikz[baseline=(char.base)]{
\node[shape=circle, draw=NovelPurple, text=NovelPurple, line width=0.8pt, inner sep=1.6pt] (char) {\small #1};
}}

\pagestyle{fancy}
\fancyhf{}
\lhead{\textbf{Novel Prep Learning Center}}
\rhead{\textbf{AP Calculus BC Study Plan}}
\cfoot{\thepage}

\titleformat{\section}{\Large\bfseries\color{NovelPurple}}{}{0em}{}
\titleformat{\subsection}{\large\bfseries\color{NovelPurple}}{}{0em}{}

\setlist[itemize]{leftmargin=1.4em}

\begin{document}

\begin{center}
{\Huge \textbf{\color{NovelPurple}AP Calculus BC Study Plan}}\\[6pt]
{\Large \textbf{Student: Andrew}}\\[4pt]
{\large \textbf{Total Instruction Time: 55 Hours}}\\[6pt]
\textbf{Novel Prep Learning Center}
\end{center}

\vspace{0.15in}

\begin{tcolorbox}[colback=LightPurple, colframe=NovelPurple, title=\textbf{Program Goal}]
This 55-hour accelerated study plan is designed for Andrew to complete the full AP Calculus BC curriculum, including limits, derivatives, integrals, differential equations, applications of integration, parametric equations, polar coordinates, vector-valued functions, and infinite series. The goal is to help Andrew build strong conceptual understanding, develop AP-style problem-solving skills, and prepare for AP Calculus BC coursework and exam-level questions.
\end{tcolorbox}

\section*{Overall Hour Distribution}

\begin{tabularx}{\textwidth}{|c|X|c|}
\hline
\rowcolor{LightPurple}
\textbf{Phase} & \textbf{Main Content} & \textbf{Hours} \\
\hline
I & Limits and Derivatives & 18 \\
\hline
II & Applications of Derivatives & 10 \\
\hline
III & Integration and Accumulation & 14 \\
\hline
IV & Differential Equations and Applications of Integration & 5 \\
\hline
V & BC-Only Topics: Parametric, Polar, and Vectors & 6 \\
\hline
VI & Infinite Sequences and Series & 8 \\
\hline
VII & AP Exam Preparation and Final Review & 4 \\
\hline
\textbf{Total} & \textbf{Complete AP Calculus BC Curriculum} & \textbf{55} \\
\hline
\end{tabularx}

\section*{Phase I: Limits and Derivatives --- 18 Hours}

\subsection*{\circnum{1} Lesson 1: Limits and Continuity --- 2 Hours}
\begin{itemize}
    \item Understand limits graphically, numerically, and algebraically.
    \item Study one-sided limits and infinite limits.
    \item Introduce continuity at a point and on an interval.
    \item Connect limits to function behavior.
\end{itemize}

\subsection*{\circnum{2} Lesson 2: Limit Laws and Algebraic Techniques --- 2 Hours}
\begin{itemize}
    \item Apply limit laws.
    \item Evaluate limits using factoring, rationalizing, and simplification.
    \item Study the Squeeze Theorem.
    \item Introduce the Intermediate Value Theorem.
\end{itemize}

\subsection*{\circnum{3} Lesson 3: Average and Instantaneous Rate of Change --- 2 Hours}
\begin{itemize}
    \item Review average rate of change.
    \item Understand instantaneous rate of change.
    \item Compare secant lines and tangent lines.
    \item Introduce the derivative definition.
\end{itemize}

\subsection*{\circnum{4} Lesson 4: Basic Derivative Rules --- 2 Hours}
\begin{itemize}
    \item Power Rule.
    \item Constant Rule.
    \item Sum and Difference Rules.
    \item Constant Multiple Rule.
    \item Higher-order derivatives.
\end{itemize}

\subsection*{\circnum{5} Lesson 5: Product Rule and Quotient Rule --- 2 Hours}
\begin{itemize}
    \item Apply the Product Rule.
    \item Apply the Quotient Rule.
    \item Practice mixed derivative problems.
    \item Emphasize correct notation and simplification.
\end{itemize}

\subsection*{\circnum{6} Lesson 6: Chain Rule --- 2 Hours}
\begin{itemize}
    \item Understand composite functions.
    \item Apply the Chain Rule.
    \item Practice multi-layer chain rule problems.
    \item Connect chain rule to AP-style derivative questions.
\end{itemize}

\subsection*{\circnum{7} Lesson 7: Trigonometric Derivatives --- 2 Hours}
\begin{itemize}
    \item Derivatives of sine, cosine, and tangent.
    \item Derivatives of secant, cosecant, and cotangent.
    \item Derivatives of inverse trigonometric functions.
    \item Mixed trigonometric derivative practice.
\end{itemize}

\subsection*{\circnum{8} Lesson 8: Exponential and Logarithmic Derivatives --- 2 Hours}
\begin{itemize}
    \item Derivatives of exponential functions.
    \item Derivatives of logarithmic functions.
    \item Natural logarithm derivative rules.
    \item Logarithmic differentiation.
\end{itemize}

\subsection*{\circnum{9} Lesson 9: Implicit Differentiation --- 2 Hours}
\begin{itemize}
    \item Understand implicit relationships.
    \item Apply implicit differentiation.
    \item Find tangent lines using implicit differentiation.
    \item Introduce higher-order implicit derivatives.
\end{itemize}

\begin{tcolorbox}[colback=SoftGray, colframe=NovelPurple, title=\textbf{Checkpoint 1}]
Quiz on limits, continuity, derivative definitions, derivative rules, chain rule, trigonometric derivatives, exponential/logarithmic derivatives, and implicit differentiation.
\end{tcolorbox}

\section*{Phase II: Applications of Derivatives --- 10 Hours}

\subsection*{\circnum{10} Lesson 10: Motion Problems --- 2 Hours}
\begin{itemize}
    \item Position, velocity, and acceleration.
    \item Interpret derivatives in motion contexts.
    \item Determine speed, direction, and acceleration.
    \item AP-style motion FRQ practice.
\end{itemize}

\subsection*{\circnum{11} Lesson 11: Increasing and Decreasing Functions --- 2 Hours}
\begin{itemize}
    \item Critical points.
    \item Increasing and decreasing intervals.
    \item First Derivative Test.
    \item Relative extrema.
\end{itemize}

\subsection*{\circnum{12} Lesson 12: Concavity and Second Derivative Test --- 2 Hours}
\begin{itemize}
    \item Concavity.
    \item Points of inflection.
    \item Second Derivative Test.
    \item Connections among $f$, $f'$, and $f''$.
\end{itemize}

\subsection*{\circnum{13} Lesson 13: Optimization --- 2 Hours}
\begin{itemize}
    \item Set up optimization equations.
    \item Identify constraints.
    \item Use derivatives to find maximum and minimum values.
    \item Practice real-world optimization problems.
\end{itemize}

\subsection*{\circnum{14} Lesson 14: Related Rates, Linearization, and L'Hospital's Rule --- 2 Hours}
\begin{itemize}
    \item Related rates setup and solving strategies.
    \item Tangent line approximation.
    \item Linearization.
    \item L'Hospital's Rule for indeterminate forms.
\end{itemize}

\begin{tcolorbox}[colback=SoftGray, colframe=NovelPurple, title=\textbf{Checkpoint 2}]
Quiz on motion, extrema, concavity, optimization, related rates, linear approximation, and L'Hospital's Rule.
\end{tcolorbox}

\section*{Phase III: Integration and Accumulation --- 14 Hours}

\subsection*{\circnum{15} Lesson 15: Riemann Sums and Area Approximation --- 2 Hours}
\begin{itemize}
    \item Left, right, and midpoint Riemann sums.
    \item Trapezoidal approximation.
    \item Interpret area as accumulation.
    \item Connect sums to definite integrals.
\end{itemize}

\subsection*{\circnum{16} Lesson 16: Definite Integrals and FTC --- 2 Hours}
\begin{itemize}
    \item Definite integral notation.
    \item Fundamental Theorem of Calculus Part I.
    \item Fundamental Theorem of Calculus Part II.
    \item Accumulation functions.
\end{itemize}

\subsection*{\circnum{17} Lesson 17: Antiderivatives and Indefinite Integrals --- 2 Hours}
\begin{itemize}
    \item Basic antiderivative rules.
    \item Indefinite integrals.
    \item Initial value problems.
    \item Interpret constants of integration.
\end{itemize}

\subsection*{\circnum{18} Lesson 18: Integration by Substitution --- 2 Hours}
\begin{itemize}
    \item Identify substitution patterns.
    \item Apply $u$-substitution.
    \item Change bounds for definite integrals.
    \item Mixed substitution practice.
\end{itemize}

\subsection*{\circnum{19} Lesson 19: Integration by Parts --- 2 Hours}
\begin{itemize}
    \item Understand the integration by parts formula.
    \item Choose $u$ and $dv$.
    \item Apply integration by parts to logarithmic and exponential functions.
    \item AP-style practice.
\end{itemize}

\subsection*{\circnum{20} Lesson 20: Partial Fractions --- 2 Hours}
\begin{itemize}
    \item Factor denominators.
    \item Set up partial fraction decomposition.
    \item Integrate rational expressions.
    \item Practice BC-style integration problems.
\end{itemize}

\subsection*{\circnum{21} Lesson 21: Improper Integrals and Mixed Integration --- 2 Hours}
\begin{itemize}
    \item Understand improper integrals.
    \item Evaluate infinite interval integrals.
    \item Evaluate discontinuous integrals.
    \item Select appropriate integration techniques.
\end{itemize}

\begin{tcolorbox}[colback=SoftGray, colframe=NovelPurple, title=\textbf{Checkpoint 3}]
Quiz on Riemann sums, definite integrals, FTC, antiderivatives, substitution, integration by parts, partial fractions, and improper integrals.
\end{tcolorbox}

\section*{Phase IV: Differential Equations and Applications of Integration --- 5 Hours}

\subsection*{\circnum{22} Lesson 22: Differential Equations and Slope Fields --- 2 Hours}
\begin{itemize}
    \item Model situations with differential equations.
    \item Sketch and interpret slope fields.
    \item Verify solutions.
    \item Apply Euler's Method.
\end{itemize}

\subsection*{\circnum{23} Lesson 23: Separation of Variables and Logistic Growth --- 1.5 Hours}
\begin{itemize}
    \item Solve separable differential equations.
    \item Apply initial conditions.
    \item Study exponential growth and decay.
    \item Introduce logistic growth models.
\end{itemize}

\subsection*{\circnum{24} Lesson 24: Applications of Integration --- 1.5 Hours}
\begin{itemize}
    \item Average value of a function.
    \item Net change and accumulation.
    \item Motion problems using integrals.
    \item Area between curves.
\end{itemize}

\section*{Phase V: BC-Only Topics --- 6 Hours}

\subsection*{\circnum{25} Lesson 25: Parametric Equations --- 2 Hours}
\begin{itemize}
    \item Define parametric equations.
    \item Differentiate parametric equations.
    \item Find second derivatives.
    \item Analyze motion using parametric equations.
\end{itemize}

\subsection*{\circnum{26} Lesson 26: Polar Coordinates --- 2 Hours}
\begin{itemize}
    \item Understand polar coordinates.
    \item Convert between polar and rectangular forms.
    \item Differentiate polar functions.
    \item Find area of polar regions.
\end{itemize}

\subsection*{\circnum{27} Lesson 27: Vector-Valued Functions and Arc Length --- 2 Hours}
\begin{itemize}
    \item Understand vector-valued functions.
    \item Differentiate and integrate vector-valued functions.
    \item Analyze velocity and acceleration.
    \item Find arc length for parametric curves.
\end{itemize}

\begin{tcolorbox}[colback=SoftGray, colframe=NovelPurple, title=\textbf{Checkpoint 4}]
Quiz on parametric equations, polar coordinates, vector-valued functions, motion, polar area, and arc length.
\end{tcolorbox}

\section*{Phase VI: Infinite Sequences and Series --- 8 Hours}

\subsection*{\circnum{28} Lesson 28: Sequences and Geometric Series --- 2 Hours}
\begin{itemize}
    \item Understand sequences.
    \item Determine convergence and divergence.
    \item Study geometric series.
    \item Apply the $n$th-term test.
\end{itemize}

\subsection*{\circnum{29} Lesson 29: Integral Test, $p$-Series, and Comparison Tests --- 2 Hours}
\begin{itemize}
    \item Integral Test.
    \item Harmonic series.
    \item $p$-series.
    \item Direct Comparison Test.
    \item Limit Comparison Test.
\end{itemize}

\subsection*{\circnum{30} Lesson 30: Alternating Series and Ratio Test --- 2 Hours}
\begin{itemize}
    \item Alternating Series Test.
    \item Alternating Series Error Bound.
    \item Ratio Test.
    \item Absolute and conditional convergence.
\end{itemize}

\subsection*{\circnum{31} Lesson 31: Taylor, Maclaurin, and Power Series --- 2 Hours}
\begin{itemize}
    \item Taylor polynomials.
    \item Maclaurin series.
    \item Lagrange Error Bound.
    \item Radius and interval of convergence.
    \item Represent functions as power series.
\end{itemize}

\begin{tcolorbox}[colback=SoftGray, colframe=NovelPurple, title=\textbf{Checkpoint 5}]
Quiz on sequences, series convergence tests, alternating series, ratio test, Taylor polynomials, Maclaurin series, and power series.
\end{tcolorbox}

\section*{Phase VII: AP Exam Preparation and Final Review --- 4 Hours}

\subsection*{\circnum{32} Lesson 32: AP Multiple Choice Strategy --- 2 Hours}
\begin{itemize}
    \item Review calculator and non-calculator strategies.
    \item Practice mixed AP-style multiple-choice questions.
    \item Identify common traps.
    \item Improve pacing and accuracy.
\end{itemize}

\subsection*{\circnum{33} Lesson 33: AP Free Response Strategy and Final Review --- 2 Hours}
\begin{itemize}
    \item Practice AP Calculus BC FRQs.
    \item Review scoring rubrics.
    \item Emphasize notation, justification, and complete reasoning.
    \item Final full-course review.
\end{itemize}

\section*{Homework and Practice Expectations}

\begin{itemize}
    \item Complete assigned practice problems after every lesson.
    \item Review all class examples before starting homework.
    \item Keep a formula and mistake notebook.
    \item Practice AP-style multiple-choice and free-response questions regularly.
    \item Revisit difficult topics weekly to strengthen long-term retention.
\end{itemize}

\section*{Assessment Plan}

\begin{tabularx}{\textwidth}{|c|X|c|}
\hline
\rowcolor{LightPurple}
\textbf{Assessment} & \textbf{Content} & \textbf{Suggested Time} \\
\hline
Checkpoint 1 & Limits and Derivatives & 45--60 min \\
\hline
Checkpoint 2 & Applications of Derivatives & 45--60 min \\
\hline
Checkpoint 3 & Integration and Accumulation & 60 min \\
\hline
Checkpoint 4 & Parametric, Polar, and Vectors & 45--60 min \\
\hline
Checkpoint 5 & Infinite Series & 60 min \\
\hline
Final Review & Mixed AP Calculus BC Review & 90--120 min \\
\hline
\end{tabularx}

\section*{Expected Outcome}

By the end of this 55-hour program, Andrew should be able to:
\begin{itemize}
    \item Understand the full AP Calculus BC curriculum.
    \item Solve standard AP-style multiple-choice questions.
    \item Write complete and organized free-response solutions.
    \item Connect graphical, numerical, analytical, and verbal representations.
    \item Use correct calculus notation and justify answers clearly.
    \item Enter AP Calculus BC with confidence and strong preparation.
\end{itemize}

\begin{tcolorbox}[colback=LightPurple, colframe=NovelPurple, title=\textbf{Teacher Note}]
This is an accelerated AP Calculus BC plan. Because the total instruction time is limited to 55 hours, Andrew should actively review after every lesson and complete homework consistently. The most important goal is not only to learn formulas, but also to understand the relationships among limits, derivatives, integrals, differential equations, and series.
\end{tcolorbox}

\vfill

\begin{center}
\textbf{\color{NovelPurple}Novel Prep Learning Center}\\
\textit{AP Calculus BC Accelerated Study Plan}
\end{center}

\end{document}